\subsubsection*{3.1 Schnelle Berechnung von $a^k \pmod{m}$}


Man schreibt im Restklassenring $\Zm{m}$, wie im Körper $\R$
\[
% KORRIGIERT: Anzahl der Faktoren und Bereich von k fehlten
  a^k := \underbrace{a \cdot a \cdot \ldots \cdot a}_{k \text{ Faktoren}} \qquad (a \neq 0,\ k \in \N,\ k \geq 1)
\]
Wie immer, wenn es um Multiplikation geht, schließen wir die 0 aus.

\begin{algorithmus}[Schnelles Potenzieren]
Nutzen die Binärdarstellung von
\[
  k = k_0 \cdot 2^0 + k_1 \cdot 2^1 + k_2 \cdot 2^2 + \ldots + k_r \cdot 2^r , \qquad k_i \in \{0, 1\}
\]
\end{algorithmus}

\begin{beispiel}
Gesucht $12^{100} \pmod{34}$
\begin{align*}
  & 100 = 2^6 + 2^5 + 2^2 \\
  \Rightarrow\quad & 12^{100} = 12^{64} \cdot 12^{32} \cdot 12^{4} \\
  & 12 \\
% KORRIGIERT: Original "12^2 \equiv 136 + 8", Rest 8 fehlte
  & 12^2 = 144 = 136 + 8 \equiv 8 \pmod{34} \\
  & 12^4 \equiv 8^2 = 64 \equiv -4 \pmod{34} \\
  & 12^8 \equiv (-4)^2 = 16 \pmod{34} \\
  & 12^{16} \equiv 16^2 = 272 - 16 \equiv -16 \pmod{34} \\
  & 12^{32} \equiv (-16)^2 \equiv -16 \pmod{34} \\
  & 12^{64} \equiv (-16)^2 \equiv -16 \pmod{34}
\end{align*}
Also
\[
  12^{100} \equiv (-16) \cdot (-16) \cdot (-4) \equiv (-16) \cdot (-4) \equiv 64 \equiv -4 \equiv 30 \pmod{34}
\]
\end{beispiel}
